Let's see the general case with k any fixed integer number. \vspace{7pt}

$\text{AtMostK}([x_1, x_2, x_3, x_4], 2)$, so we would like to set up an encoding that allows us to check if the maximum number of bit set to $1$ is exactly 
$2$, using two Full Adders. As before, we distinguish boolean decision variables and the constraint related to them:
\begin{itemize}
    \renewcommand{\labelitemi}{-}
    \item $a_0$, $b_0$, $a_1$, $b_1$ take values of $x_1$, $x_2$, $x_3$ and $x_4$.
    \item We want at most two boolean variables set to $1$, therefore the constraint is:
    $$\neg c_0 \wedge (c_2 \rightarrow (\neg d_0 \wedge \neg d_1 \wedge \neg c_1))$$ 
    We are just performing the sum between binary numbers: starting from the notion $2_2 = 10$, we can use this information in order to describe our constraint.
    In this way, we know for sure that $c_2$ cannot be assigned to true and only $d_0$ or $d_1$ can be assigned to $1$, providing that the final result 
    does not exceed the fixed integer number.
\end{itemize}
However, considering $x_1 + x_2 + x_3 \leq 0$ we assume that by unit propagation the solver will set all these variables to false, without any additional 
operation from its list. As usual, we can exploit an encoding using one Full Adder:
\begin{itemize}
    \renewcommand{\labelitemi}{-}
    \item $a_0$, $b_0$ and $c_0$ take values of $x_1$, $x_2$ and $x_3$.
    \item In addition, our constraint is divided into:
    \begin{enumerate}
        \item $\neg d_0 \wedge \neg c_1$
        \item $d_0 \leftrightarrow (x_1 \wedge \neg x_2 \wedge \neg x_3) \vee (\neg x_1 \wedge x_2 \wedge \neg x_3) \vee (\neg x_1 \wedge \neg x_2 \wedge x_3) \vee (x_1 \wedge x_2 \wedge x_3)$
        \item $c_1 \leftrightarrow (x_1 \wedge x_2) \vee (x_1 \wedge x_3) \vee (x_2 \wedge x_3)$
    \end{enumerate}
    \item Converting them into CNF form the whole thing becomes:
    \begin{enumerate}[start = 2]
        \item $\neg d_0 \rightarrow (\neg x_1 \wedge x_2 \wedge x_3) \vee (x_1 \wedge \neg x_2 \wedge x_3) \vee (x_1 \wedge x_2 \wedge \neg x_3) \vee (\neg x_1 \wedge \neg x_2 \wedge \neg x_3)$
        \item $\neg c_1 \leftrightarrow (\neg x_1 \wedge \neg x_2) \vee (\neg x_1 \wedge \neg x_3) \vee (\neg x_2 \wedge \neg x_3)$
    \end{enumerate}
\end{itemize}
Unit propagation cannot propagate anything, since we have more than one free variable in both cases, for instance:
\begin{itemize}
    \renewcommand{\labelitemi}{-}
    \item $\neg d_0 \rightarrow (\neg x_1 \wedge x_2 \wedge x_3) \equiv \neg d_0 \vee (\neg x_1 \wedge x_2 \wedge x_3)$, $3$ different free variables, which one should we propagate?
    \item $\neg c_1 \leftrightarrow (\neg x_1 \wedge \neg x_2) \equiv \neg c_1 \vee (\neg x_1 \wedge \neg x_2)$, $2$ different free variables, which one should we propagate?
\end{itemize}
Therefore, without any additional operation, using exclusively unit propagation, our solver cannot derive that any boolean variable should be set to false. 
Adder encoding is not arc-consistent encoding (remember the two key points described in \ref{s_7_9})!